\section{Did the shortlist contain the school the pupil later attended?}\label{sec:observed-search}
The predictors were trained on offers. We can test them against a different outcome recorded months later by linking each original applicant to the national school register for the following April 30 \citep{registerdata}. Of 2,044,649 applicant-years, 1,969,120 link to a registered school. A menu scores one hit if any displayed course is at that school. The target is the same historically attended school for both rules. For the placement comparison, we still clear the entire cohort, then count placements among the register-linked applicants.

\begin{table}[ht]\centering\small
\caption{Historical school enrollment predicted and placements in replay. The middle columns count linked applicants whose five-choice menu, including any protected option, contains their April school. The last column is the simulated placement change for those same applicants.}\label{tab:enrolled-target}
\begin{tabular}{@{}lrrrrr@{}}\toprule
Application year & Linked & Blind hits & Informed hits & Hit change & Placement change\\\midrule
2022 & 549,997 & 444,227 & 445,234 & $+1,007$ & $-405$\\
2023 & 516,921 & 420,820 & 422,018 & $+1,198$ & $-482$\\
2024 & 457,231 & 374,906 & 376,291 & $+1,385$ & $-434$\\
2025 & 444,971 & 365,988 & 367,036 & $+1,048$ & $-288$\\\midrule
Total & 1,969,120 & 1,605,941 & 1,610,579 & $+4,638$ & $-1,609$\\\bottomrule
\end{tabular}
\end{table}

The informed rule gained 6,184 hits and lost 1,546 relative to the blind rule (Table~\ref{tab:enrolled-target}). Requiring the school, type of education and grade to match gives a net gain of 4,487. Could missing registers reverse the school-level result? For an unlinked pupil, the blind rule can gain a hit over the informed rule only if it lists a school absent from the informed menu. Just 3,760 unlinked records meet that condition. The reverse condition holds for 3,757. Allowing each of the 75,529 unlinked records any school or no enrollment gives a full-cohort net gain between $4,638-3,760=878$ and $4,638+3,757=8,395$. The lower bound is positive in each year (54, 313, 442 and 69). These are bounds from the two menus, with no assumption about why a register link is missing.

The average conceals a limit precisely where families needed another option. Among the 126,708 linked pupils still without a regular offer after the wait list, only 25,499 (20.12\%) attended a school anywhere on their eligible original report. The blind and informed menus contained the eventual school for 25,158 and 25,122, respectively. Even a perfect predictor restricted to that report would miss nearly four in five. What happened to these families after their first search?

\paragraph{Where families searched next}
A submitted list is a record of one search. To see whether families subsequently pursued other schools, we linked the regular and complementary application files by masked pupil identifier within each of the four years. The complementary round allows eligible families to submit another list from schools with vacancies \citep{complementary}. We selected everyone whose initial and post-wait-list regular outcomes both recorded no offer. We counted a new school only when its identifier was absent from \emph{all} original entries and the family explicitly listed a course there in a same-grade complementary application. Automatic continuity entries and assignments to a course the family never listed do not count as new explicit offers. Finally we linked the initial outcome in the complementary round to the next school year's April 30 enrollment register by pupil and school \citep{saedata,registerdata}.

\begin{table}[ht]\centering\small
\caption{Observed choices after the regular round, in application records. Each column is a subset of the column to its left. Enrollment is at the new school where the complementary round offered the listed course.}\label{tab:actual-search}
\begin{tabular}{@{}lrrrrr@{}}\toprule
Year & No regular offer & Reapplied$^{a}$ & Listed new school & Offered it & Enrolled there\\\midrule
2022 & 43,077 & 21,371 & 20,858 & 15,095 & 8,192\\
2023 & 37,757 & 18,537 & 18,116 & 13,895 & 7,923\\
2024 & 30,486 & 14,599 & 14,532 & 12,340 & 7,304\\
2025 & 30,025 & 16,041 & 15,730 & 13,071 & 7,280\\\midrule
Total & 141,345 & 70,548 & 69,236 & 54,401 & 30,699\\\bottomrule
\end{tabular}\\[2pt]
{\footnotesize $^{a}$Same reported grade in both rounds. Another 595 records returned at a different grade.}
\end{table}

Of the 141,345 initially and finally unoffered regular applicants, 21.72\% subsequently enrolled at a school absent from the original list through this observed sequence. Among 54,401 offered an explicitly added school, 51,956 appeared anywhere in the following April's register and 30,699 appeared at that school (56.43\% of all 54,401). Another 21,257 appeared at a different school and 2,445 had no register link. Among the 101,209 linked pupils who finished the regular wait list without an offer and whose April school was outside their eligible original report, the complementary offers in Table~\ref{tab:actual-search} account for 30,699 (30.33\%). Search often expanded even when a former option remained on the public vacancy list: 45,848 of 46,940 such same-grade returners added a school. New supply, new information and an unsuccessful first round can all change a family's decision. The linkage establishes when an earlier omission preceded later enrollment. It cannot identify which family would pursue a particular omitted school under a hypothetical five-choice recommendation.

This distinction matters for the comparison with our simulated losses. Of the 2,527 applicants placed under the blind predictor alone, 509 went without an offer through the actual regular wait list, and 103 of those 509 later enrolled at an explicitly added school through the sequence above. Much of the simulated loss population thus started the actual second round in a different position. Applying the 21.72\% rate to predicted losses would mix populations and interventions.

A separate Chilean field study helps motivate which responses to examine. Near the threshold for a nonplacement-risk warning, \citet{smartplatforms} estimated a 21.6 percentage point increase in the probability of adding any school (SE 1.0), a 20.5 point increase in adding at the end (SE 0.9), and a 3.8 point increase in placement at a listed choice (SE 0.9). We verified the published estimates against the authors' archived output tables \citep{warningdata}. Those families received advice about risk. The simulations below ask how our two proposed menus would compare if applicants recovered options, including when they pursued an additional option last.
